% ============================================================================
% Appendix E: Validation, Ablation, and Deployment Guide
% Sources: experiments_v1/appendix/ sections B2, D1, F3
% ============================================================================

\section{Validation Methodology, Ablation Studies, and Deployment Guide}
\label{app:validation_ablation}

% ----------------------------------------------------------------------------
% E.1: Statistical Validation (from B2)
% ----------------------------------------------------------------------------

\subsection{Statistical Validation of Semantic Structure}
\label{app:validation_methodology}

We perform comprehensive validation to ensure our Alignment Tax findings (Figure~\ref{fig:alignment_tax}) are statistically robust and not artifacts of methodological choices.

\subsubsection{Statistical Significance Testing}

We validate the cluster distinction using multiple tests robust to distributional assumptions:

\paragraph{Primary Tests.}
\begin{itemize}[leftmargin=*,noitemsep,topsep=0pt]
    \item \textbf{Mann-Whitney $U$-test (non-parametric):} $p < 2.86 \times 10^{-143}$
    \item \textbf{Welch's $t$-test (parametric):} $p < 2.36 \times 10^{-92}$
    \item \textbf{Effect size (Cohen's $d$):} $1.901$ (large effect, $d > 0.8$)
    \item \textbf{Confidence intervals (95\%):} Low PC1: $[+0.113, +0.153]$, High PC1: $[-0.738, -0.625]$ (non-overlapping)
\end{itemize}

\paragraph{Distribution Diagnostics.}
Shapiro-Wilk normality tests indicate both distributions are non-normal ($p < 0.001$), justifying non-parametric tests. Both parametric and non-parametric tests converge on highly significant differences.

\subsubsection{Threshold Validation}

To ensure the $\text{PC1} = 0.3$ decision boundary is principled rather than arbitrary, we validated it through multiple independent methods:

\paragraph{Grid Search.}
We evaluated 50 candidate thresholds across the PC1 range using a composite score combining silhouette coefficient, Davies-Bouldin index, mean reward gap separation, and cluster balance. The optimal threshold by composite score is $\text{PC1} = 0.317$, with $\text{PC1} = 0.3$ achieving nearly identical performance (composite score difference: 0.2\%). \textbf{Multiple testing note:} Because the threshold was selected from 50 candidates, the primary significance tests are subject to threshold selection bias. After Bonferroni correction for 50 comparisons, the Mann-Whitney $U$-test remains highly significant ($p_{\text{corrected}} < 10^{-141}$). We further validate that results are robust to threshold perturbation: all thresholds in $[0.2, 0.4]$ yield $p < 10^{-100}$ (see Sensitivity Analysis below).

\paragraph{Unsupervised Clustering.}
Independent clustering methods identify consistent boundaries:
\begin{itemize}[leftmargin=*,noitemsep,topsep=0pt]
    \item $k$-means clustering ($k=2$): boundary $\approx 0.148$
    \item Gaussian Mixture Model ($k=2$): boundary $\approx 0.462$
    \item Mean across methods: $0.320 \pm 0.105$
\end{itemize}
The threshold $\text{PC1} = 0.3$ falls within one standard deviation of the cross-method mean.

\paragraph{Sensitivity Analysis.}
Results remain highly significant across threshold perturbations: all thresholds in $[0.2, 0.4]$ achieve $p < 10^{-100}$ with strong silhouette scores ($> 0.39$).

\subsubsection{High-Dimensional Structure Validation}

To verify that the bimodal structure is not a dimensionality reduction artifact, we examined cluster quality across multiple dimensionalities:

\begin{table}[h]
\centering
\caption{Cluster Quality Across Dimensionalities}
\label{tab:multidim_validation}
\begin{tabular}{lccc}
\toprule
\textbf{Space} & \textbf{Dimensions} & \textbf{Silhouette} & \textbf{Separation Ratio} \\
\midrule
2D PCA & 2 & 0.495 & --- \\
32D PCA & 32 & 0.255 & 1.41 \\
384D Embeddings & 384 & 0.057 & 0.81 \\
\bottomrule
\end{tabular}
\end{table}

Silhouette scores decrease in higher dimensions due to the curse of dimensionality. However, the cluster assignment based on PC1 remains predictive of reward gaps in the original 384D space (Spearman $\rho = -0.395$, $p < 10^{-70}$), confirming that PC1 captures task-relevant semantic structure. The 5.4\% variance captured by PC1+PC2 represents the \emph{task-relevant} semantic axis; the remaining 94.6\% encompasses task-orthogonal variation (topic, language, style) that does not affect model performance trade-offs.

\subsubsection{Data Quality Analysis}

The dataset contains zero exact duplicates (100\% unique: 1,871/1,871). Near-duplicate analysis (cosine similarity $\geq 0.95$) identifies only 201 pairs among 330 High PC1 prompts (0.37\% rate). Intra-cluster diversity analysis shows High PC1 diversity score of 0.355 and Low PC1 diversity score of 0.953, confirming findings are not dominated by specific templates.

% ----------------------------------------------------------------------------
% E.2: Corralling Ablation (from D1)
% ----------------------------------------------------------------------------

\subsection{Corralling Hyperparameter Ablation}
\label{app:corralling_ablation}

To characterize sensitivity to hyperparameter choices, we conducted a systematic grid search over 15 configurations spanning 5 learning rates ($\eta \in \{0.1, 0.5, 1.0, 2.0, 5.0\}$) and 3 mixing parameters ($\gamma \in \{0.0, 0.05, 0.10\}$), with 3 independent seeds per configuration (45 total experiments). All experiments use $N=1{,}121$ labeled samples and 3 models (Mixtral, GPT-4-Turbo, GPT-4o).

\begin{table}[h]
\centering
\caption{Corralling hyperparameter ablation. Mean $\pm$ std across 3 seeds, ranked by cumulative regret.}
\label{tab:corralling_ablation}
\small
\begin{tabular}{cccccc}
\toprule
$\eta$ & $\gamma$ & Seeds & Regret (mean $\pm$ std) & Reward (mean $\pm$ std) & Rank \\
\midrule
\textbf{5.0} & \textbf{0.10} & 3 & \textbf{59.33 $\pm$ 3.40} & \textbf{0.9390 $\pm$ 0.0030} & \textbf{1st} \\
1.0 & 0.00 & 3 & 59.67 $\pm$ 10.27 & 0.9387 $\pm$ 0.0092 & 2nd \\
5.0 & 0.05 & 3 & 61.67 $\pm$ 2.05 & 0.9370 $\pm$ 0.0018 & 3rd \\
0.1 & 0.05 & 3 & 64.67 $\pm$ 3.86 & 0.9343 $\pm$ 0.0034 & 4th \\
0.1 & 0.10 & 3 & 68.33 $\pm$ 1.70 & 0.9310 $\pm$ 0.0015 & 5th \\
1.0 & 0.10 & 3 & 70.67 $\pm$ 9.74 & 0.9289 $\pm$ 0.0087 & 6th \\
2.0 & 0.00 & 3 & 75.00 $\pm$ 22.76 & 0.9251 $\pm$ 0.0203 & 7th \\
0.5 & 0.10 & 3 & 85.33 $\pm$ 27.45 & 0.9158 $\pm$ 0.0245 & 8th \\
0.5 & 0.05 & 3 & 86.00 $\pm$ 27.82 & 0.9153 $\pm$ 0.0248 & 9th \\
2.0 & 0.10 & 3 & 86.33 $\pm$ 28.80 & 0.9150 $\pm$ 0.0257 & 10th \\
0.5 & 0.00 & 3 & 88.33 $\pm$ 31.63 & 0.9132 $\pm$ 0.0282 & 11th \\
1.0 & 0.05 & 3 & 90.00 $\pm$ 48.19 & 0.9117 $\pm$ 0.0430 & 12th \\
0.1 & 0.00 & 3 & 94.33 $\pm$ 41.68 & 0.9078 $\pm$ 0.0372 & 13th \\
5.0 & 0.00 & 3 & 94.33 $\pm$ 18.70 & 0.9078 $\pm$ 0.0167 & 14th \\
2.0 & 0.05 & 3 & 99.00 $\pm$ 54.46 & 0.9037 $\pm$ 0.0486 & 15th \\
\bottomrule
\end{tabular}
\end{table}

\paragraph{Key Findings.}
\begin{enumerate}[nosep]
    \item \textbf{Learning Rate:} Higher rates ($\eta \geq 1.0$) enable faster expert adaptation. The optimal $\eta=5.0$ achieves aggressive weight updates while maintaining stability through the exploration floor ($\gamma=0.10$).
    
    \item \textbf{Exploration Floor:} The mixing parameter ($\gamma > 0$) prevents expert starvation and stabilizes learning. Configurations with $\gamma=0$ exhibit $14\times$ higher variance (e.g., $\eta=2.0$, $\gamma=0.0$: std=54.46 vs std=3.40 for optimal), while $\gamma=0.10$ yields consistent performance.
    
    \item \textbf{Reproducibility:} The optimal configuration demonstrates low variance across seeds (std=3.40), confirming results are not due to fortunate initialization.
    
    \item \textbf{Robustness:} Performance remains competitive across $\eta \in [0.5, 5.0]$ when $\gamma \geq 0.05$, indicating reasonable tolerance to hyperparameter misspecification.
\end{enumerate}

\paragraph{Convergence Validation.}
Log-log regression (excluding 50-sample burn-in) yields growth exponent $\beta=0.669 \pm 0.002$ with $R^2=0.9903$, confirming sublinear regret growth $O(T^{0.669})$ and satisfying the PAC-learnability criterion ($\beta < 1$).

\paragraph{Semantic Transfer for Cold-Start Models.}
Since warmup priors cover only Mixtral and GPT-4-Turbo, GPT-4o is initialized via semantic transfer from GPT-4-Turbo with conservative scaling ($\gamma=0.05$). The bandit rapidly discovers GPT-4o as the dominant choice (70.8\% usage), correctly identifying the warmup expert's GPT-4-Turbo bias (79.9\% in isolation) as suboptimal.

% ----------------------------------------------------------------------------
% E.3: Deployment Strategy Guide (from F3)
% ----------------------------------------------------------------------------

\subsection{Deployment Strategy Selection Guide}
\label{app:strategy_guide}

Table~\ref{tab:strategy_guide} provides a practical decision framework for practitioners selecting a routing strategy based on prior quality assessment.

\begin{table}[h]
\centering
\caption{%
    \textbf{Deployment Strategy Selection Guide.}
    Regret values from holdout evaluation 
    (N=750 prompts, severe domain mismatch: 68.6\%$\to$13.7\% hard prompts).
    *Performance in matched scenario expected to exceed Corralling.
}
\label{tab:strategy_guide}
\small
\begin{tabular}{@{}llccp{4.0cm}@{}}
\toprule
\textbf{Prior Quality} & \textbf{Strategy} & \textbf{Regret} & \textbf{vs Best} & \textbf{When to Use} \\
\midrule
Unknown / Uncertain
    & Corralling & 59.2 $\pm$ 7.1 & +20\% & 
    Cross-domain deployment, risk-averse, prior quality unvalidated \\
\midrule
Known Bad
    & Tabula Rasa & \textbf{49.5 $\pm$ 2.8} & \textbf{--} & 
    Cold start, severe mismatch, validated $<$50\% accuracy \\
\midrule
Known Good
    & Warmup Only & N/A* & N/A & 
    Same domain as warmup, validated $>$80\% accuracy \\
\bottomrule
\end{tabular}
\end{table}

\paragraph{Validation Protocol.}
Evaluate warmup predictions on $N=100$--$200$ deployment samples. Accuracy $>80\%$ = good priors (use Warmup Only); $50$--$80\%$ = uncertain (use Corralling); $<50\%$ = bad priors (use Tabula Rasa).

\paragraph{Key Findings.}
(1) Under severe mismatch (as in the original ablation), Corralling hedges between experts, incurring modest overhead when priors are transferable;
(2) Tabula Rasa is optimal when priors are known-bad;
(3) Strategy selection depends on prior quality assessment (see Validation Protocol above).
All configurations use constant $\alpha=2.0$ (validated in Appendix~\ref{sec:alpha_ablation}).

